Sind Injektionen \(F\colon A\inj B\) und \(G\colon B\inj A\)
gegeben, so gibt es eine Bijektion zwischen \(A\) und \(B\).
Die folgenden Aussagen beschreiben die dazu verwendete Teilmenge und
die beiden bijektiven Zweige.

Die vollständigen Beweise dieses Abschnitts stehen in zwei verknüpften
Begleitfassungen: Die \CSBReadingLink{} entwickelt die Konstruktion
in einem zusammenhängenden mathematischen Beweis; die \CSBProofsLink{}
führen die einzelnen Ableitungsschritte aus. Beide verwenden die
hier angegebenen Satznummern.

\Needspace{14\baselineskip}
\CSBPartDef
\Needspace{14\baselineskip}
\CSBPartSubset
\Needspace{14\baselineskip}
\CSBPartContainsSeed
\Needspace{14\baselineskip}
\CSBPartMinimal
\Needspace{16\baselineskip}
\CSBPartClosed
\begingroup
  \ProofPartSetParent{\theformulaThm}
  \setcounter{proofpartnr}{0}
  \CSBAuxReference{CantorBernsteinPartClosedUniversal}
\endgroup
\Needspace{14\baselineskip}
\CSBPartFixedPoint
\Needspace{16\baselineskip}
\CSBComplementIdentity
\begingroup
  \ProofPartSetParent{\theformulaThm}
  \setcounter{proofpartnr}{0}
  \CSBAuxReference{CantorBernsteinComplementForward}
  \CSBAuxReference{CantorBernsteinComplementBackward}
\endgroup
\Needspace{14\baselineskip}
\CSBSecondBranchImage
\Needspace{14\baselineskip}
\CSBFixedInjections
\noindent\emph{Beweis:} \CSBReadingLink{} und \CSBProofsLink{}.
